\documentclass[10pt]{article}
\usepackage[T1]{fontenc}
\usepackage{lmodern,microtype,amsmath,amssymb,array,tabularx,booktabs}
\usepackage[margin=0.85in]{geometry}
\usepackage{xurl,hyperref}
\hypersetup{colorlinks=true,urlcolor=blue,pdftitle={Appendix A review -- Version 18}}
\setlength{\parindent}{0pt}
\setlength{\parskip}{5pt}
\newcolumntype{P}[1]{>{\raggedright\arraybackslash}p{#1}}
\newcolumntype{Y}{>{\raggedright\arraybackslash}X}
\begin{document}
{\LARGE Appendix A: range comparison and literature review}\par
{\large Version 18, 31 October 2026}\par

Version 18 adds a numbered comparison of regular images and ranges,
uses it in the three applications, and supplies literature precedents
beside the corresponding constructions. The standalone appendix is
Appendix A; its mathematical body is Appendix B of the manuscript.
The six-section organization and the accepted manuscript abstract are retained.

\section*{Current outline}
\renewcommand{\arraystretch}{1.15}
\begin{tabularx}{\textwidth}{@{}P{0.075\textwidth}P{0.295\textwidth}Y@{}}
\toprule
Section & Title & Function and dependencies\\\midrule
A.1 & Categorical hypotheses and notation & Regular category, stable disjoint
coproducts for syntax, Boolean operations and conditionals, decidable variable
equality, signature objects, and parameter stability.\\
A.2 & Terms, formulas, and free variables & Initial syntax and list algebras,
constructor morphisms, free-variable morphisms, and the sentence subobject;
uses A.1.\\
A.3 & Admissible substitution & Parameterized structural recursion, total term
substitution, formula substitution on a complemented admissibility subobject,
and identity and absent-variable laws; uses A.1--A.2.\\
A.4 & Inference rules and deductively closed subobjects & A slice endofunctor
for inference rules and Theorem A.5: the regular image of an initial derivation
algebra is the least closed subobject. Requires regularity and finite coproducts.\\
A.5 & First-order theories & Logical schemes, formula and sentence theories,
Theorem A.6, and Lemma A.7 identifying an admissible regular image with a range;
uses A.2--A.4.\\
A.6 & Propositional theories & The propositional instance of A.4, reusing the
three propositional schemes of A.5.\\\bottomrule
\end{tabularx}

The central construction takes an axiom subobject to the least sentence theory
containing it. Regular images enter the closure theorem and its pullback to
sentences. The earlier constructor and substitution constructions use initial
algebras, finite limits, and the specified coproduct structure.

\section*{Range comparison and applications}
Lemma A.7 (B.7 in the manuscript) assumes a split restriction category
$\mathcal S$ with ranges, a regular total category
$\mathcal D=\operatorname{Tot}(\mathcal S)$, and a regular-image factorization
$f=me$ of a total morphism. If $m:I\hookrightarrow Y$ has a partial inverse
$r:Y\rightharpoonup I$, then
\[
\operatorname{ran}(f)=mr=\bar r.
\]
The proof shows that $mr$ receives $f$ and is below every receiving restriction
idempotent. Splitting such an idempotent gives a total factorization of $f$;
the universal property of its regular image supplies the required comparison.
The agreement law is the stated range convention, not an extra conclusion of
this proof. The partial-map application cites the construction of split ranges
from a stable factorization system and the agreement condition for an
epimorphic left class.

Sections 4 and 8 now cite B.7 for the first-order theorem ranges, using their
existing regularity and admissible-monomorphism results. Section 9 cites B.7
for the propositional theorem range in its partial-map category.
The final paragraph of A.5 records that a restriction functor preserving ranges
sends the displayed splitting to a splitting of the image range. Numerical
realization proofs remain in the applications.

\newpage
\section*{Literature and exact roles}
The search covered many-sorted initial syntax, categorical variable binding,
parameterized list recursion, finitary deduction, factorization systems on
algebras, and range restriction categories. The references below were checked
for their particular role; no priority claim is made.

\begin{tabularx}{\textwidth}{@{}P{0.31\textwidth}Y@{}}
\toprule
Source and location & Relationship to the appendix\\\midrule
Goguen--Thatcher--Wagner--Wright (1977), \S2 & Initial many-sorted algebras
for abstract syntax. This is an antecedent, not an existence theorem for the
appendix's arbitrary ambient category.\\
Goguen--Burstall (1986), \S4 & Initial-algebra treatment of term and sentence
sorts in logical languages.\\
Cockett (1990); Maietti (2010), Definition 2.4 and following paragraph &
Parameterized list recursion and its formulation by pullback into slices.
The attribution to Cockett was checked through Maietti's primary text and
bibliographic records; Cockett's full text was unavailable.\\
Fiore--Plotkin--Turi (1999), \S\S2--4 & Initial-algebra accounts of binding
and substitution, with bound-variable renaming identified.\\
Gabbay--Pitts (2002), \S6 and proof of Corollary 6.7 & Binding syntax modulo
bound-variable renaming; the product construction retaining the original
syntax component obtains structural recursion from initial-algebra iteration.\\
Katsumata (2004), \S2, Example 1(3) & An explicit first-order example with
quantification in a functor-category treatment of binding syntax.\\
Aczel (1977), Propositions 1.1.4 and 1.4.5 & Finitary inductive closure as
conclusions of finite proofs. The appendix expresses rule data in a slice.\\
Wi\ss{}mann (2022), Lemma 3.6 & Factorization systems lift to algebras when
the endofunctor preserves the left class. Theorem A.5 includes its own proof
of the needed regular-epimorphism descent.\\
Cockett--Guo--Hofstra (2012), I: Theorems 4.5, 4.7, Lemma 4.12 and \S4.11;
II: \S3.16, Proposition 3.17 & Split ranges and factorization systems,
the stronger agreement law, and the additional conditions relating these
factorizations to regular images.\\\bottomrule
\end{tabularx}

The appendix keeps raw named-variable syntax. Its formula substitution is
restricted to the admissibility subobject. This difference from the cited
binding constructions is stated where those constructions are discussed;
their existence results are not transferred to arbitrary $\mathcal D$.

The existing references for restriction categories and the logical axiom
schemes remain. Full bibliographic entries and links appear in the appendix.
Fiore--Hur's second-order equational logic was considered but not used.
Maietti--Vickers's induction principle for consequence in arithmetic universes
was examined; its specific induction theorem is not the closure-as-image
theorem proved here. Neither was added as a dependency. A reference to Makkai
was not added merely for general relevance.

\newpage
\section*{Writing and notation review}
The governing prose source was Knuth, Larrabee, and Roberts,
\emph{Mathematical Writing}, Stanford report STAN-CS-89-1193 (1989),
\href{https://www-cs-faculty.stanford.edu/~knuth/papers/cs1193.pdf}{official text}.
Section 1 was consulted for connected exposition, definitions at first use,
consistent notation, and readable sentences. Citation relevance and exact
categorical typing were also checked as editorial and mathematical requirements.

\begin{tabularx}{\textwidth}{@{}P{0.17\textwidth}P{0.29\textwidth}Y@{}}
\toprule
Location & Finding or issue & Implemented action and rationale\\\midrule
A.2 opening & ``Constructor objects'' had no separate definition. &
Replaced by ``signature objects, syntax objects, and constructor morphisms'';
uses the declared types (\S1, points 11 and 14).\\
A.2 comparison & Unqualified ``substitution'' would also include total term
substitution. & Specified ``formula substitution''; avoids ambiguity
(\S1, point 17).\\
A.3 recursion & The structural-recursion construction has a direct precedent. &
Cited Gabbay--Pitts at the product construction and preserved the proof.
This is a source-attribution decision.\\
A.5 range lemma & The comparison needs regularity and an admissible image
inclusion. & Both are in the lemma statement; the proof types every comparison
morphism. This is a mathematical requirement.\\
A.5 partial maps & The lemma's global range hypothesis needs justification in
this instance. & Added the split-range construction and its literature
dependency before applying the lemma.\\
A.1 and A.3 & Product powers, truth injections, and conditional branches were
already specified. & Retained the version 17 conventions and copairing;
no notation redesign was needed.\\\bottomrule
\end{tabularx}

The definition-order check confirmed that $V$ belongs to $\mathcal D$ and has
decidable equality before its use; $K$ and its action on morphisms precede the
syntax definition; admissibility precedes the logical schemes; and ranges,
restriction monomorphisms, and the split total category precede Lemma A.7.
No reorganization of the six sections was needed.

\section*{Verification and scope}
Separate targeted checks covered initial syntax and list precedents, binding
precedents and bibliographic data, deduction precedents, the range comparison,
and its three applications. Final read-only checks of the new passages led to
the wording and partial-map-hypothesis repairs above. The complete revised
appendix source was read for continuity and editorial-process language.

The standalone appendix and assembled manuscript compiled without unresolved
references, citation warnings, or box warnings. All standalone pages and the
changed application and literature pages in the manuscript were rendered and
visually inspected. The standalone and integrated appendix bodies agree apart
from publication metadata and the integrated compatibility label. Every new
download has a version 18 filename; PDF metadata and displayed versions agree.

This was a targeted mathematical and literature review, not a fresh review of
every theorem in the paper or an exhaustive priority search. The earlier
four-agent hostile-referee, writing, outline, and reconciliation reports remain
at \href{https://github.com/flengyel/CategoricalEssentialIncompleteness/tree/main/Audit/2026-10-31-appendix-v14}{the version 14 review directory}.
No new Lean verification was performed.
\end{document}
